% sections/introduction.tex
% Note: this paper uses \sfrac{a}{b} from xfrac, NOT \nicefrac{a}{b} (nicefrac.sty unavailable).
\section{Introduction}
\label{sec:intro}

Modern deep reinforcement learning for continuous control rests on a
substantial collection of algorithmic patches. Proximal Policy
Optimization~\cite{schulman2017ppo} adds a clipped surrogate, value
baseline, generalised advantage estimation, observation
normalisation, reward scaling, and entropy regularisation to the
basic policy gradient~\cite{williams1992}. Soft
Actor-Critic~\cite{haarnoja2018sac} adds two target networks, a
replay buffer, automatic temperature tuning, and twin Q-functions.
Each addition has been justified by an empirical problem -- variance,
off-policyness, exploration -- but together they form a system whose
individual contributions are difficult to isolate.

Two careful studies have argued that much of the gain attributed to
specific algorithms is in fact attributable to implementation
details. Engstrom et~al.~\cite{engstrom2020} found that code-level
optimisations explain most of PPO's lead over TRPO.
Andrychowicz et~al.~\cite{andrychowicz2021} trained nearly a quarter
of a million agents and identified roughly fifty hyperparameters
whose interactions dominate performance. Both papers conclude that
the standard tricks matter; we ask the converse:
\emph{which tricks could be removed if the underlying objective were
better behaved?}

We propose a minimalist baseline, \textsc{MinPG}, built from three
ideas:
\begin{enumerate}[leftmargin=*]
  \item A truncated Student-$t$ policy, with state-dependent degrees of
    freedom, replaces the standard Tanh-Gaussian. Heavy tails and a
    full $C^\infty$ density act as a mollifier that smooths the
    contact-discontinuous viability margin.
  \item A soft-minimum over per-step margins replaces the task return
    as the \emph{sole} training signal. There is no shaped task
    reward, no exploration bonus, no curriculum-aware multiplier.
  \item A Frobenius-norm capacity clamp on every layer replaces $L_2$
    weight decay with a hard Lipschitz bound, ensuring bounded score
    variance without a learned baseline.
\end{enumerate}
Given these three pieces, the score-function gradient is sufficient.
We use no value function, no replay buffer, no PPO clipping, no
entropy bonus, no DAgger correction. The simulator is treated as a
black box; no gradient flows through physics.

\paragraph{Contributions.}
(1)~We give a smoothness lemma showing that the Student-$t$ mollifier
restores $C^\infty$ dependence of the soft-min expected objective on
$\theta$, even when the underlying margins are non-differentiable
through contact transitions (Section~\ref{sec:theory}).
(2)~We prove a finite-variance bound for the score-function estimator
under the capacity clamp, eliminating the standard variance argument
for value baselines.
(3)~We benchmark \textsc{MinPG} against PPO and SAC on Walker2d-v4
(Section~\ref{sec:exp}). \emph{[Empirical claim placeholder pending
experiments.]}
(4)~Ablations isolate the contribution of each component:
Student-$t$ vs.\ Gaussian, soft-min vs.\ task return, with vs.\
without capacity clamp.

\paragraph{Scope.}
This is a study of \emph{whether} the minimalist combination is
viable on a single locomotion task, not a claim of state-of-the-art
across continuous control. The paper succeeds if \textsc{MinPG} is
competitive with default-configured PPO and SAC on Walker2d-v4 given
the margin formulation; it fails as an empirical claim if the
variance of REINFORCE swamps the mollifier benefit. We hold
exploration mechanisms, representation learning, and sim-to-real
transfer constant: this paper isolates a single axis of variation.
